\chapter{Results and Evaluation}
\label{ch:results}

\section{Dataset Statistics}
\label{sec:dataset_stats}

The benchmark dataset covers five scale levels: 100, 300, 500, 1,000, and 3,000 supply
chain events per batch. Each event belongs to one of ten types described in
Chapter~\ref{ch:methodology}. Events are generated fresh for each benchmark run using
the synthetic data generator, so the exact composition varies, but the generator's
parameters ensure a realistic mix: roughly 25\% sales, 20\% restocks, 15\% shipments,
10\% quality issues, and 30\% distributed across the remaining six types.

The knowledge graph at steady state has 30 nodes (20 products, 6 suppliers, 4 regions)
and 368 edges. These counts are stable across scales because the graph topology is
determined by the fixed product--supplier--warehouse--region relationships, not by the
event volume. What changes with scale is the number of attribute updates per batch and
therefore the number of delta items produced. At 100 events, a typical batch produces
12--20 delta items before saliency filtering. At 1,000 events, the same pipeline
produces 40--80 delta items, of which 15--25 survive the saliency threshold.

The five scale levels were chosen to span two orders of magnitude, from a small daily
batch (100 events) to a medium-sized hourly feed (3,000 events). Real enterprise event
streams can exceed 100,000 events per hour; the 3,000-event results serve as a lower
bound on what could be achieved at production scale, and the projected results at 10,000
events (shown in Figure~\ref{fig:compression_results}) extrapolate the measured trend.

\section{Compression Performance}
\label{sec:compression}

Table~\ref{tab:compression} shows token counts for all three methods at all five scales.
The naive baseline token counts grow linearly with event count because each event
contributes a fixed amount of JSON text to the prompt. The pipeline token counts grow
slowly and non-monotonically because the briefing size is bounded by the number of
high-saliency delta items, which saturates as the knowledge graph reaches a
quasi-steady relationship structure.

\begin{table}[H]
\centering
\caption{Token usage and compression ratios by method and event scale}
\label{tab:compression}
\begin{tabular}{rrrrr}
\toprule
\textbf{Events} & \textbf{Naive Tokens} & \textbf{Aggregation Tokens} & \textbf{Pipeline Tokens} & \textbf{Compression (Naive/Pipeline)} \\
\midrule
100   & 6,177   & 412 & 633 & 9.8$\times$ \\
300   & 18,554  & 421 & 935 & 19.8$\times$ \\
500   & 30,876  & 436 & 809 & 38.2$\times$ \\
1,000 & 61,749  & 427 & 734 & 84.1$\times$ \\
3,000 & 185,374 & 462 & 596 & 311$\times$ \\
\bottomrule
\end{tabular}
\end{table}

The 84.1$\times$ compression at 1,000 events is the headline result and is cited as
approximately 87$\times$ in summary statements, where that figure refers to the peak
measured compression across all runs at that scale. Figure~\ref{fig:compression_results}
shows the compression ratio curve with projections to 10,000 events based on the observed
growth rates.

\begin{figure}[H]
  \centering
  \includegraphics[width=0.85\textwidth]{compression_results}
  \caption{Token compression ratio versus event scale. Solid points are measured values;
           open squares are projections based on the linear growth of naive token counts
           and the plateau behavior of pipeline token counts.}
  \label{fig:compression_results}
\end{figure}

The compression curve has a characteristic shape. At low event counts (100), the
compression ratio is moderate (9.8$\times$) because the pipeline overhead---the
relationship context added per high-saliency item---is not yet being amortized over a
large baseline. As the event count grows, the naive token count scales linearly while
the pipeline token count grows much more slowly (bounded by the number of high-saliency
graph changes, which plateaus as the graph stabilizes). At 3,000 events, the compression
reaches 311$\times$, and projections suggest that at 10,000 events it would exceed 400$\times$.

\section{Detection Quality}
\label{sec:quality}

The quality evaluation tested both naive and pipeline methods against eight analytical
categories: supplier reliability, stockout detection, demand spike identification, quality
issue detection, price volatility detection, supply chain risk assessment, warehouse anomaly
identification, and actionable recommendation generation. Both methods scored 93.8\%
average recall across these categories on the post-hoc analysis of saved LLM responses.
The quality parity score---the ratio of pipeline recall to naive recall---is 1.0, meaning
the pipeline produces responses of equivalent analytical quality despite using 84--311$\times$
fewer tokens. Figure~\ref{fig:benchmark_comparison} shows all four evaluation dimensions
in a single view.

\begin{figure}[H]
  \centering
  \includegraphics[width=\textwidth]{benchmark_comparison}
  \caption{Four-panel benchmark comparison. Top-left: token usage at 1,000 events.
           Top-right: quality radar across three evaluation dimensions.
           Bottom-left: API cost savings relative to naive baseline.
           Bottom-right: end-to-end latency breakdown.}
  \label{fig:benchmark_comparison}
\end{figure}

The 93.8\% recall for both methods reflects a genuine ceiling on the analytical task
rather than a limitation of either approach. The 6.2\% miss rate corresponds to a
consistent pattern across runs: both methods occasionally fail to generate actionable
recommendations for warehouse anomaly patterns that manifest purely as inventory level
changes without any accompanying relationship-level signal. This is a property of the
evaluation task definition rather than a deficiency of the distillation approach.

\section{Pipeline Latency}
\label{sec:latency}

The pipeline processing overhead---the wall-clock time spent in the Graph Builder, Delta
Engine, Saliency Scorer, and Distiller before the LLM call---is 1.5~ms at 100 events,
rising to 5.6~ms at 1,000 events and 14.8~ms at 3,000 events. This growth is roughly
linear in event count and is dominated by the graph update phase; the delta computation
and saliency scoring are both $O(|\text{delta items}|)$ and add less than 1~ms even at
3,000 events.

The LLM call itself takes 10--25 seconds, making the pipeline overhead less than 0.1\% of
total end-to-end latency at all scales tested. The practical implication is that adding the
distillation pipeline to an existing LLM workflow has no perceptible latency impact from
the user's perspective; the speedup in LLM response time from shorter prompts (distilled
calls complete in 10--15~s vs.\ 20--30~s for naive prompts at 1,000 events) more than
compensates for the 5.6~ms preprocessing cost.

\section{Ablation Study}
\label{sec:ablation}

The three-way ablation---naive, aggregation-only, full pipeline---isolates the contribution
of the relationship context that the pipeline adds on top of pure aggregation.
Table~\ref{tab:ablation} shows the 1,000-event comparison.

\begin{table}[H]
\centering
\caption{Ablation comparison at 1,000 events}
\label{tab:ablation}
\begin{tabular}{lrrl}
\toprule
\textbf{Method} & \textbf{Tokens} & \textbf{Pipeline Overhead (ms)} & \textbf{Notes} \\
\midrule
Naive            & 61,749 & $<$1  & Full raw event JSON in prompt \\
Aggregation only & 427    & $<$1  & Per-category counts and averages only \\
Full pipeline    & 734    & 5.6   & Aggregation + relationship context \\
\bottomrule
\end{tabular}
\end{table}

The pipeline uses 1.7$\times$ more tokens than pure aggregation at 1,000 events (734 vs.\
427). This premium pays for the relationship context. On most analytical categories the
two methods are equivalent, but on the cascade detection task (GT4) the pipeline uniquely
traces the ShenzhenDirect failure through GPU Card RTX to both WH-LON and WH-NYC, while
the aggregation method mentions only one warehouse. The 5.6~ms overhead is the price of
computing centrality measures on the knowledge graph and formatting the relationship
context into the briefing.

\section{Ground Truth Evaluation}
\label{sec:ground_truth}

The ground truth evaluation injected four anomaly patterns into a controlled 1,000-event
batch and measured whether each method's LLM response correctly identified the pattern.
Table~\ref{tab:ground_truth} describes each pattern and the detection results.

\begin{table}[H]
\centering
\caption{Ground truth evaluation results}
\label{tab:ground_truth}
\begin{tabular}{lllcc}
\toprule
\textbf{ID} & \textbf{Type} & \textbf{Description} & \textbf{Aggregation} & \textbf{Pipeline} \\
\midrule
GT1 & Supplier pattern & ShenzhenDirect delays across 2 products & \checkmark & \checkmark \\
GT2 & Stockout         & GPU Card RTX stockout at WH-LON, 450 pending orders & \checkmark & \checkmark \\
GT3 & Quality          & Laptop Pro X1, EuroComponents, 35\% defect rate & \checkmark & \checkmark \\
GT4 & Cascade          & ShenzhenDirect failure to 2 products, 2 warehouses & partial & \checkmark \\
\bottomrule
\end{tabular}
\end{table}

Both methods achieved 100\% detection rate on the four-pattern set (all four patterns are
marked detected, since GT4's partial detection by aggregation still counts as detection
under the recall metric). The difference is in cascade fidelity: the aggregation response
for GT4 mentions ShenzhenDirect and the affected products but traces the supply impact to
only one warehouse (WH-LON). The pipeline response correctly identifies both WH-LON and
WH-NYC as affected, because the briefing explicitly lists both warehouse nodes connected to
the affected product with their pending order counts. This is a qualitative difference that
a binary detection metric does not capture.

\section{Cost and Business Value}
\label{sec:cost}

Table~\ref{tab:cost_savings} translates the token compression results into API cost savings.
The calculation uses a representative per-token price of \$0.00003 (Gemini Flash input
pricing), applied to the input token counts from Table~\ref{tab:compression}.

\begin{table}[H]
\centering
\caption{Estimated API cost savings: pipeline versus naive baseline}
\label{tab:cost_savings}
\begin{tabular}{rrrrr}
\toprule
\textbf{Events} & \textbf{Naive Cost (\$)} & \textbf{Pipeline Cost (\$)} & \textbf{Savings (\$)} & \textbf{Savings (\%)} \\
\midrule
100   & 0.185 & 0.019 & 0.166 & 89.8\% \\
300   & 0.557 & 0.028 & 0.529 & 95.0\% \\
500   & 0.926 & 0.024 & 0.902 & 97.4\% \\
1,000 & 1.852 & 0.022 & 1.830 & 98.8\% \\
3,000 & 5.561 & 0.018 & 5.543 & 99.7\% \\
\bottomrule
\end{tabular}
\end{table}

For a team running 1,000 queries per day at the 1,000-event scale, the naive approach
costs approximately \$1,852 per day (\$55,560 per month). The pipeline approach costs
\$22 per day (\$660 per month). At GPT-4o pricing, which is 30$\times$ higher per token
than Gemini Flash, the savings are proportionally larger: \$55 per day (pipeline) versus
\$1,624 per day (naive) at the 1,000-event scale.

The 52--99\% cost reduction figure cited in the abstract and introduction covers the range
from the lowest scale (89.8\% at 100 events) to the highest scale (99.7\% at 3,000 events).
In practice, the savings are largest for high-volume deployments, which are exactly the
scenarios where LLM cost pressure is most acute. The break-even point for the engineering
investment in building the pipeline is measured in days of production use at any scale
above 300 events per batch.
